\section{Experimental Setup}
\label{sec:evaluation-experimental-setup}

In this section, I will describe the experimental setup and the methods used to evaluate the performance of the various optimizations and compare them against the baseline implementation and analytical results, but before that I want to go over the test systems and the relevant parameters and how they impacted the result.

\subsection{Test Systems}
\label{sec:evaluation-experimental-setup-test-systems}

Testing was done both locally on my own machine, and on specific nodes of the TU Berlin HPC. 
The specifications of both systems can be seen in Table~\ref{tab:test-systems}. 
As already mentioned in \ref{sec:background-cuda-gpu-architecture}, FP64 performance on gaming/consumer such as the RTX 3080 is neutered when compared to the A100, despite both chips being part of the same architecture generation.
On the other hand, the Intel Xeon E5-2630v4 of the A100 clusters is quite old, and is handily beaten by the R7 5800X3D in both single-core and multi-core performance, which makes comparing performance a little awkward.
% To compare both systems fairly I needed to benchmark the pure C++ and CUDA codes separately.
% didn't end up doing it...

\begin{table}[htbp]
    \centering
    \resizebox{\textwidth}{!}{%
    \begin{tabular}{lccccc}
        \toprule
         & Local & gpu[022,025,026] & gpu[023,024,027] & gpu068 & gpu[069-072] \\
        \midrule
        CPU & AMD R7 5800X3D & \multicolumn{2}{c}{Intel Xeon E5-2630v4} & \multicolumn{2}{c}{AMD EPYC 9555} \\
        RAM & 32 GB DDR4 & \multicolumn{2}{c}{512 GB DDR4} & \multicolumn{2}{c}{1.1 TB DDR5} \\
        GPU & RTX 3080 & Tesla A100 & 2xTesla A100 & 4xH200 NVL & 2xH200 NVL \\
        VRAM & 10GB & 40GB & 80GB & 564GB & 282GB \\
        FP32:FP64 & 64:1 & \multicolumn{2}{c}{2:1} & \multicolumn{2}{c}{2:1} \\
        Driver / CUDA & 580.173.02 / 13.0 & \multicolumn{2}{c}{570.86.10 / 12.8} & \multicolumn{2}{c}{---} \\
        \bottomrule
    \end{tabular}%
    }
    \caption{Specifications of the systems used for testing.}
    \label{tab:test-systems}
\end{table}

Locally, I ran NGPB using the \snakecase{mpirun -n <ranks>} \snakecase{ngpb --prmfile options.prm} command.
The TU HPC uses Slurm, an open-source scalable cluster management and job scheduling system for large and small Linux clusters \cite{slurm_overview}. 
I used Claude used it to create a Singularity recipe \cite{kurtzer2017singularity} to build a portable container that includes all the necessary dependencies and NGPB itself.
This container could then run on the cluster without the need to separately install everything, which of course would not be allowed to begin with.
I built the image locally using \texttt{apptainer build --fakeroot}, and copied it to the cluster using \texttt{scp} together with the Slurm script and the necessary input files, where I ran it using the \texttt{sbatch} command.
I also generated small scripts to launch NGPB on a single cluster node, which all generally followed the same scheme as the example shown in Listing~\ref{lst:run-slurm}.
Launching on multiple nodes was more difficult, since NGPB uses a different OpenMPI implementation than the cluster.
Instead of running with \snakecase{mpirun} from inside the container, I had to use \snakecase{srun --mpi=pmix} to spawn on container per rank.
I measured the communication latency on the TU HPC, can be seen in Table~\ref{tab:mpi-latency-table}.
Logically, I chose Unified Communication X (UCX) for the transport layer, though it still is a bit slower than \snakecase{mpirun} on a single container.

% Generated by: python3 scripts/mpi_latency_table.py test5/mpidiag2-1883214.out
% Pass --cases A,B,E,G,H,J to drop the intermediate diagnostic transports.
% ------------------------------------------------------------------------
%  MPI transport measurements on the TU Berlin HPC.
%  Generated by scripts/mpi_latency_table.py from test5/mpidiag2-1883214.out -- regenerate
%  rather than edit. Include with % ------------------------------------------------------------------------
%  MPI transport measurements on the TU Berlin HPC.
%  Generated by scripts/mpi_latency_table.py from test5/mpidiag2-1883214.out -- regenerate
%  rather than edit. Include with % ------------------------------------------------------------------------
%  MPI transport measurements on the TU Berlin HPC.
%  Generated by scripts/mpi_latency_table.py from test5/mpidiag2-1883214.out -- regenerate
%  rather than edit. Include with \input{figures/mpi_latency_table}
%  Needs \usepackage{booktabs}.
% ------------------------------------------------------------------------
\begin{table}[htbp]
    \centering
    \footnotesize
    \begin{tabular}{llrrr}
        \toprule
        placement & transport & ping-pong & bandwidth & allreduce \\
                  &           & [$\mu$s]  & [GB/s]    & [$\mu$s]  \\
        \midrule
        1 node, 2 ranks & \texttt{mpirun}, one container & 0.25 & 8.73 & 0.38 \\
        1 node, 2 ranks & \texttt{srun}+pmix, default & \multicolumn{3}{c}{deadlocks} \\
        1 node, 2 ranks & vader, emulated single-copy & 0.26 & 5.20 & 0.36 \\
        1 node, 2 ranks & \texttt{btl=tcp,self} & 5.17 & 4.49 & 7.08 \\
        1 node, 2 ranks & \texttt{pml=ucx} & 0.37 & 7.00 & 0.39 \\
        1 node, 2 ranks & \texttt{pml=ucx}, no shared memory & 3.16 & 4.91 & 2.64 \\
        \midrule
        2 nodes, 1 rank each & default & 48.42 & 0.31 & 52.60 \\
        2 nodes, 1 rank each & \texttt{pml=ucx} & 6.10 & 0.60 & 5.73 \\
        2 nodes, 1 rank each & \texttt{btl=tcp,self} & 33.73 & 0.34 & 50.27 \\
        \midrule
        2 nodes, 2 ranks each & \texttt{pml=ucx} & 0.31 & 7.94 & 7.41 (4 ranks) \\
        \bottomrule
    \end{tabular}
    \caption{MPI performance between two nodes of the TU Berlin HPC, measured
        with \texttt{mpi\_bench.c} from inside the NGPB container. Ping-pong and
        allreduce use 8\,B messages, bandwidth 8\,MB. Ping-pong is always rank 0
        to rank 1, which is a same-node pair in the last row.}
    \label{tab:mpi-latency-table}
\end{table}

%  Needs \usepackage{booktabs}.
% ------------------------------------------------------------------------
\begin{table}[htbp]
    \centering
    \footnotesize
    \begin{tabular}{llrrr}
        \toprule
        placement & transport & ping-pong & bandwidth & allreduce \\
                  &           & [$\mu$s]  & [GB/s]    & [$\mu$s]  \\
        \midrule
        1 node, 2 ranks & \texttt{mpirun}, one container & 0.25 & 8.73 & 0.38 \\
        1 node, 2 ranks & \texttt{srun}+pmix, default & \multicolumn{3}{c}{deadlocks} \\
        1 node, 2 ranks & vader, emulated single-copy & 0.26 & 5.20 & 0.36 \\
        1 node, 2 ranks & \texttt{btl=tcp,self} & 5.17 & 4.49 & 7.08 \\
        1 node, 2 ranks & \texttt{pml=ucx} & 0.37 & 7.00 & 0.39 \\
        1 node, 2 ranks & \texttt{pml=ucx}, no shared memory & 3.16 & 4.91 & 2.64 \\
        \midrule
        2 nodes, 1 rank each & default & 48.42 & 0.31 & 52.60 \\
        2 nodes, 1 rank each & \texttt{pml=ucx} & 6.10 & 0.60 & 5.73 \\
        2 nodes, 1 rank each & \texttt{btl=tcp,self} & 33.73 & 0.34 & 50.27 \\
        \midrule
        2 nodes, 2 ranks each & \texttt{pml=ucx} & 0.31 & 7.94 & 7.41 (4 ranks) \\
        \bottomrule
    \end{tabular}
    \caption{MPI performance between two nodes of the TU Berlin HPC, measured
        with \texttt{mpi\_bench.c} from inside the NGPB container. Ping-pong and
        allreduce use 8\,B messages, bandwidth 8\,MB. Ping-pong is always rank 0
        to rank 1, which is a same-node pair in the last row.}
    \label{tab:mpi-latency-table}
\end{table}

%  Needs \usepackage{booktabs}.
% ------------------------------------------------------------------------
\begin{table}[htbp]
    \centering
    \footnotesize
    \begin{tabular}{llrrr}
        \toprule
        placement & transport & ping-pong & bandwidth & allreduce \\
                  &           & [$\mu$s]  & [GB/s]    & [$\mu$s]  \\
        \midrule
        1 node, 2 ranks & \texttt{mpirun}, one container & 0.25 & 8.73 & 0.38 \\
        1 node, 2 ranks & \texttt{srun}+pmix, default & \multicolumn{3}{c}{deadlocks} \\
        1 node, 2 ranks & vader, emulated single-copy & 0.26 & 5.20 & 0.36 \\
        1 node, 2 ranks & \texttt{btl=tcp,self} & 5.17 & 4.49 & 7.08 \\
        1 node, 2 ranks & \texttt{pml=ucx} & 0.37 & 7.00 & 0.39 \\
        1 node, 2 ranks & \texttt{pml=ucx}, no shared memory & 3.16 & 4.91 & 2.64 \\
        \midrule
        2 nodes, 1 rank each & default & 48.42 & 0.31 & 52.60 \\
        2 nodes, 1 rank each & \texttt{pml=ucx} & 6.10 & 0.60 & 5.73 \\
        2 nodes, 1 rank each & \texttt{btl=tcp,self} & 33.73 & 0.34 & 50.27 \\
        \midrule
        2 nodes, 2 ranks each & \texttt{pml=ucx} & 0.31 & 7.94 & 7.41 (4 ranks) \\
        \bottomrule
    \end{tabular}
    \caption{MPI performance between two nodes of the TU Berlin HPC, measured
        with \texttt{mpi\_bench.c} from inside the NGPB container. Ping-pong and
        allreduce use 8\,B messages, bandwidth 8\,MB. Ping-pong is always rank 0
        to rank 1, which is a same-node pair in the last row.}
    \label{tab:mpi-latency-table}
\end{table}


\lstinputlisting[style=bashstyle, caption={Example Slurm submission script (\texttt{run.slurm}) used to launch NGPB on the TU Berlin HPC.}, label={lst:run-slurm}]{../test5/run.slurm}

\subsection{NGPB Parameters}
\label{sec:evaluation-experimental-setup-ngpb-parameters}

% the third paragraph can stay here as it explains why I picked which parameter values
In the NGPB paper \cite{diflorio2025nextgenpb}, the authors showed that a \texttt{perfil1} of 90\%-95\% or 0.9-0.95 is perfectly sufficient, and would reduce the runtime and memory footprint significantly compared to 0.8.
Out of curiosity, I did run a couple of tests using a \texttt{perfil1} of 1. It didn't change the result all that much, and it also didn't make the tests run any faster, so I stuck with the author-approved values.
\texttt{perfil2} is kept at its default value of 0.2, as the coarse region has a much lower impact on performance. 
\texttt{scale} is also kept at its default value of 2, which was proven to be sufficient in the paper.
To my surprise, \snakecase{loc_refinement} yielded worse results in my testing, and served only to test \snakecase{energy_cuda}.

When it came to the choice between the two solvers I settled on LIS pretty quickly.
A short test on my local machine showed that mumps takes way longer than LIS for 1CCM with all the settings being the same aside from the solver.
Running the same test with 6VYB with \texttt{perfil1} set to 0.95 LIS managed to finish, while mumps froze.
That does not mean mumps is entirely useless. My guess it that it is meant to run on huge clusters, not on a single machine.
On the other hand this demonstrates again how much more flexible linear solvers are compared to direct solvers.

\subsection{MPI Rank}
\label{sec:evaluation-experimental-setup-mpi-ranks}

The number of ranks for MPI is set through the \texttt{-n} or \texttt{-np} flags. 
As already explained in \ref{sec:nextgenpb-libraries-bimpp}, the total memory consumption of all nodes increases slightly with the number of ranks, as each one needs to keep track of ghost- and non-local elements.
% omitting the "measured" values as I don't think they're a good metric
However, it can also slightly influence the result itself, as the individual results of the ranks may be reduced in a different order.
The magnitude of this error is at FP64 noise level, but can impact smaller molecules like 1CCM, where I discovered that the flux charge varied wildly with \texttt{-n}, though this behavior is unfortunately to be expected.
With large molecules like 6VYB and 1VSZ this error was imperceptible.
Most local testing was performed using an \texttt{-n} of 4, while testing on the cluster was performed using an \texttt{-n} of 16.

\subsection{Compiler Flags}
\label{sec:evaluation-experimental-setup-compiler-flags}

Aside from the additional compiler flags added due to inclusion of CUDA, I largely used the same flags from the original makefile and \snakecase{local_settings}. 
Although initially I changed the optimization flag from \texttt{-O2} to \texttt{-O0} due to fears of catastrophic cancellations during floating-point calculations, spawned by the flux charge variations described in \ref{sec:evaluation-experimental-setup-mpi-ranks}, it quickly proved more trouble than it was worth. 
It made the code extremely slow, which exaggerated the speedup achieved by the GPU and the perceived need to optimize other parts of the code. 
Besides that, the dependencies of NGPB were also compiled using \texttt{-O2}. 

I reverted to \texttt{-O2}, but purposefully stayed away from the commented-out \texttt{-Ofast -mtune=native -march=native} setting lingering inside the \snakecase{local_settings}. While it would have provided even faster baseline performance, \texttt{-Ofast} in particular implies \texttt{-ffast-math}, which permits reassociation of floating point sums, assumes finite arithmetic, and on x86 links \texttt{crtfastmath.o} that enables flush-to-zero and denormals-are-zero process-wide, thereby affecting the linked dependencies as well.
This is a notable relaxation of the IEEE semantics and could compromise the CPU baseline against which the GPU was compared.
The other two flags \texttt{-mtune=native} and \texttt{-march=native} are more benign. They change instruction scheduling and enable FMA contraction, the latter being enabled by default on \texttt{nvcc}.

\subsection{Testing}
\label{sec:evaluation-experimental-setup-testing}

The main two metrics used for evaluation were the execution time and accuracy of the results, while memory usage served a secondary role. 
I created a handful of test cases using different PQR models as seen in Table~\ref{tab:test-molecules}.
Since 1CCM took only a few seconds on the local machine, I did not use it on the cluster, as the speedup would have been too hard to determine there.
Conversely I could not run H1N1 locally, as it needed hundreds of gigabytes of RAM.

The test cases generally fell into two categories: analytical test cases and comparative test cases.
The analytical test cases were used to verify the correctness of the results with respect to the known analytical results taken from the NGPB paper \cite{diflorio2025nextgenpb}. 
Those allowed me to verify the error magnitudes without having to rely on the objective correctness of the baseline, and also sanity-check the parameters used to achieve the original results.
They were not used to compare speed, as they are fairly small and thus not very computationally expensive.

Comparative test cases, on the other hand, were used to compare the results of the optimized implementation against the baseline, which isn't actually formed by the results of the completely unmodified version of NGPB. 
Instead it's the result produced by the version containing the triangulation fix described in \ref{sec:implementation-energy-calculation-triangulation-fix}, as it is more accurate than the result produced by the original code.
Unlike the analytical benchmarks, these were used to compare both speed and accuracy, as their models were much larger.

Table~\ref{tab:test-cases} lists the test cases together with the NGPB configuration used for each. 
Model settings were kept identical across all of them. 
Mesh settings were mostly kept at the given values, though I already pointed out instances where I experimented with different values.

\begin{table}[htbp]
    \centering
    \footnotesize
    \begin{tabular}{p{0.9cm}p{2.8cm}p{2.0cm}p{2.8cm}p{4.7cm}}
        \toprule
        Test & PQR File & Category & Mesh Settings & Model Settings \\
        \midrule
        test0 & 1CCM.pqr & Comparative & \snakecase{mesh_shape}=0, \texttt{perfil1}=0.95,\newline\texttt{perfil2}=0.2, \texttt{scale}=2.0 &
        \multirow{8}{4.7cm}{\snakecase{bc\_type}=1\newline$\varepsilon_m$=2, $\varepsilon_s$=80\newline\snakecase{ionic\_strength}=0.145\,mol/L\newline$T$=298.15\,K\newline\snakecase{calc\_energy}=2\newline\snakecase{calc\_coulombic}=1} \\
        test1 & 6VYB.pqr & Comparative & \snakecase{mesh_shape}=0, \texttt{perfil1}=0.95,\newline\texttt{perfil2}=0.2, \texttt{scale}=2.0 & \\
        test2 & 1vsz.pqr & Comparative & \snakecase{mesh_shape}=0, \texttt{perfil1}=0.95,\newline\texttt{perfil2}=0.2, \texttt{scale}=2.0 & \\
        test3 & UNITSPHERE\_-\newline DISC\_2.pqr & Analytical & \snakecase{mesh_shape}=0, \texttt{perfil1}=0.15,\newline\texttt{perfil2}=0.15, \texttt{scale}=2.0 & \\
        test4 & 30sfere.pqr & Analytical & \snakecase{mesh_shape}=0, \texttt{perfil1}=0.20,\newline\texttt{perfil2}=0.20, \texttt{scale}=2.0 & \\
        test5 & 6VYB.pqr (cluster) & Comparative & \snakecase{mesh_shape}=0, \texttt{perfil1}=0.95,\newline\texttt{perfil2}=0.2, \texttt{scale}=2.0 & \\
        test6 & 1vsz.pqr (cluster) & Comparative & \snakecase{mesh_shape}=0, \texttt{perfil1}=0.95,\newline\texttt{perfil2}=0.2, \texttt{scale}=2.0 & \\
        test7 & H1N1.pqr (cluster) & Comparative & \snakecase{mesh_shape}=0, \texttt{perfil1}=0.95,\newline\texttt{perfil2}=0.2, \texttt{scale}=0.769231 & \\
        \bottomrule
    \end{tabular}
    \caption{Test cases and their NGPB configuration parameters.}
    \label{tab:test-cases}
\end{table}

Naturally, the versions of the dependencies had to be kept static. I settled on using NanoShaper version 1.5 taken from the master branch instead of 0.8, as version 1.5 features much greater performance, without affecting accuracy.
Since the optimizations to BIMPP benefit both the original code and the CUDA-accelerated versions without affecting the accuracy, I will accept them into the baseline.
The versions of the remaining dependencies were kept according to the Dockerfile within the NGPB repository.
To greatly accelerate the capture of data, I extensively used AI generated scripts, which collected an plotted data automatically.

\subsection{Comparing performance}
\label{sec:evaluation-experimental-setup-comparing-speed}

Thankfully, NGPB already prints out the time that each stage took. 
This is done using two helper macros called TIC() and TOC(), which are portable throughout the code base. 
Those rough timings are enough to show whether an optimization is faster, since the computation time for the larger molecules is a lot longer than small timer randomness, or other run-to-run-variances.
Unless a figure or table states otherwise, every timing reported in this chapter is the median of 3 repeated runs. 1CCM's short runtime makes it more susceptible to this noise, so I instead used 10 repeats for it.

For more exact timings in the GPU-specific optimizations, I profiled the CUDA kernels using \texttt{ncu}, which runs them multiple times, eliminating runtime variances.
Additionally it provides detailed information about their utilization of GPU resources and gives suggestions on how to optimize them. 
I was able to simply wrap it around the \texttt{mpirun} command. 
This recorded the separate kernel launches of each rank, which wasn't really a problem though. 
Unfortunately, I could only use \texttt{ncu} locally since it requires \snakecase{NVreg_RestrictProfilingToAdminUsers} to be disabled.
On top of that, an FP64-intensive kernel running on the RTX 3080 might obscure other performance bottlenecks.
Due to this, I was not able to get the full picture needed to optimize the CUDA code to the fullest extent.
I avoided profiling the linear solvers running on the GPU, since I made no code changes to them, by selecting the built-in solvers like LIS instead. 
This way I could avoid putting profiler start and stop flags into the code.

\subsection{Comparing accuracy}
\label{sec:evaluation-experimental-setup-comparing-accuracy}

One of the goals of this thesis was to maintain NGPB's level of accuracy measured by the relative error magnitude of the various energies, as shown in Table~\ref{tab:ngpb-error-magnitudes}. 
The analytical reference values, seen in Table~\ref{tab:ngpb-analytical-values}, were graciously provided to me by the authors of NGPB, as the decimal places of the ones in the paper are truncated too much to be useful.
% TODO I am not sure if I should include the formulas for calculating the analytical values.

\begin{table}[htbp]
    \centering
    \begin{tabular}{lccc}
        \toprule
         & Polarization & Ionic & Total \\
        \midrule
        Kirkwood sphere & 7.38e-10 & 3.39e-02 & 1.72e-04 \\
        30-sphere       & 4.16e-05 & 1.39e-02 & 7.46e-04 \\
        \bottomrule
    \end{tabular}
    \caption{NGPB relative error magnitudes for the Kirkwood sphere and 30-sphere test systems.}
    \label{tab:ngpb-error-magnitudes}
\end{table}

\begin{table}[htbp]
    \centering
    \begin{tabular}{lccc}
        \toprule
         & Polarization & Ionic & Total \\
        \midrule
        Kirkwood sphere & -68.30597989 & -0.3480318551 & -68.65401175 \\
        30-sphere       & -10310.56623446509 & -151.12876079887 & -2254.400164369044 \\
        \bottomrule
    \end{tabular}
    \caption{Analytical reference values for the Kirkwood sphere and 30-sphere test systems.}
    \label{tab:ngpb-analytical-values}
\end{table}

The combination of my modifications to NGPB and its subsystems must not increase those magnitudes for all analytical test cases, meaning the magnitude of their error contributions must be at the same order or lower. 
On the other hand, error magnitudes smaller than or equal to $10^{-16}$ are considered FP64 machine-level noise, so aiming for those is not necessary, giving me some wiggle room between $10^{-10}$ and $10^{-16}$ for possible approximations.

Comparing the accuracy between the linear solvers should be unnecessary, since in theory all linear solvers aim to converge towards a set residual or accuracy value, which is completely independent of the solver or preconditioning.
However in practice, I noticed that LIS was somehow achieving higher accuracy than AMGX despite the same target accuracy of $10^{-6}$.
Upon closer inspection, LIS oversteps its target and achieves a residual of $10^{-8}$, while AMGX quits right after reaching it.
To rule out any degradation of the Kirkwood test result through AMGX, I pinned its accuracy to $10^{-10}$.
%Still this makes accuracy comparisons obsolete because AMGX was set to a lower tolerance.
%nah I'll actually compare accuracy
This does put AMGX at a theoretical disadvantage, since higher accuracy requires more iterations, but LIS's strange overstepping-behavior already makes this comparisons more difficult than it should have been.

The accuracy of the energy calculation can be verified independently using the comparative test cases.
Additionally, both the Kirkwood Sphere test, as well as the 30 sphere test are a poor fit for FMM, because they are too small, and their interaction list collapses fully into the near-field.
Since it is supposed to follow the math exactly and the triangulation between the baseline and GPU optimization didn't change, both results should only deviate by the aforementioned error magnitudes, as long as the intermediate results of the previous stages stayed the same.

\subsection{Comparing memory utilization}
\label{sec:evaluation-experimental-setup-comparing-memory-utilization}

I used Heaptrack to profile the RAM utilization, which wasn't the main focus during testing, as most of it was caused by BIMPP, not NGPB.
Unlike \texttt{ncu}, wrapping \texttt{heaptrack} around \texttt{mpirun} produced useless data, probably from the MPI code itself. 
Instead it had to be called between \texttt{mpirun} and \texttt{ngpb} which produced one profile for each rank. 
Since rank 0 is a de facto master rank, which performs important coordination work, it can be easily identified, as it uses more memory than the other ranks.

Register spills and cache utilization of the CUDA kernels directly tie into performance and were given higher priority.
One can be diagnosed at compile time without using \texttt{ncu} using the \texttt{--ptxas-options=-v} flag of \texttt{nvcc}.
The other is naturally part of the \texttt{ncu} profile.

Outside of the cuDSS experiment, VRAM utilization was basically dwarfed by the RAM usage, especially during the energy calculation. 
It was seen more as a correctness criterion: ``Did it fit into VRAM at all?'', especially on my local system with only 10GB.